% ECONOMICS_PROBLEM: {"id": "J16-461", "title": "Origins of child penalties", "jel_code": "J16", "source_id": "OP-0457"}
% !TeX program = xelatex
\documentclass[11pt,a4paper]{article}
\usepackage[margin=23mm]{geometry}
\usepackage{fontspec}
\setmainfont{Latin Modern Roman}
\usepackage{amsmath,amssymb,mathtools}
\usepackage{xcolor,enumitem,booktabs,longtable,array,xurl,hyperref}
\definecolor{muted}{HTML}{53636C}
\hypersetup{colorlinks=true,urlcolor=blue,linkcolor=blue}
\setlength{\parindent}{0pt}
\setlength{\parskip}{5pt}
\newcommand{\R}{\mathbb R}
\newcommand{\E}{\mathbb E}
\newcommand{\N}{\mathbb N}
\newcommand{\ind}{\mathbf 1}
\newcommand{\Var}{\operatorname{Var}}
\newcommand{\Cov}{\operatorname{Cov}}
\newcommand{\supp}{\operatorname{supp}}
\DeclareMathOperator*{\argmin}{arg\,min}
\DeclareMathOperator*{\argmax}{arg\,max}
\newcommand{\entrymeta}[3]{{\sffamily\small\color{muted}\textbf{\texttt{#1}}\quad|\quad #2\par #3\par}}
\newcommand{\Problemref}[1]{\texttt{#1}}
\newcommand{\JELref}[2]{\texttt{#2} (JEL \texttt{#1})}
\begin{document}
\section*{Origins of child penalties}
\textbf{Problem ID:} \texttt{J16-461}\quad \textbf{JEL:} \texttt{J16}\par
% BEGIN ECONOMICS STATEMENT
\label{problem:OP-0457}
{\sffamily\small\color{muted}Original subject group: Education, families and demography\par}
\label{definitions:problem:30007739}
\textbf{Audit classification:} Background; proposed extension. Defined controlled policy-by-parenthood outcomes; randomized policy alone does not identify the no-birth trajectory.
\subsection*{Definitions and precise research target}
\textit{Editorial formalization.} Consider couples expecting a first child in one country, with both parents employed before childbirth. Let \(a\in\{0,1\}\) denote a specified policy that changes childcare access or fathers’ nontransferable leave, and let \(Y_{igt}(a,1)\) be real earnings for parent \(g\) in couple \(i\) at event time \(t\), where \(g\in\{m,f\}\) indexes mothers and fathers and \(t=0\) is the first birth. Let \(Y_{igt}(a,0)\) denote earnings under a specified no-birth counterfactual. Define the gendered parenthood gap at each \(t\in\{1,\ldots,10\}\) by
\[P_t(a)=E[(Y_{imt}(a,1)-Y_{imt}(a,0))-(Y_{ift}(a,1)-Y_{ift}(a,0))].\]
Identify \(P_t(1)-P_t(0)\) using credible policy variation and a defensible no-birth trajectory, rather than attribute the entire event-study gap to discrimination or preferences. Collect employer responses, hours, childcare use, partner time, and stated expectations; distinguish mechanisms only with additional variation or explicit exclusion restrictions. Address selection into parenthood and leave eligibility, and retain parents who leave employment.

Define policy \(a\in\{0,1\}\) and a controlled parenthood regime \(b\in\{0,1\}\), where \(b=1\) fixes the first-birth date at zero and \(b=0\) delays births beyond the horizon. Let \(Y_{igt}(a,b)\) be earnings, with zeros after labor-force exit. Then \(P_t(a)=E[Y_{imt}(a,1)-Y_{imt}(a,0)-Y_{ift}(a,1)+Y_{ift}(a,0)]\) is a policy-by-parenthood interaction for the baseline cohort; no-birth policies must be explicitly meaningful (leave is unused in \(b=0\)). If policy changes birth timing, use total policy effects on a preannouncement cohort instead of conditioning on policy-selected parents. Randomized policy alone does not identify all four childbirth potential outcomes. The event-study gender gap is observational unless the no-birth trajectory and birth-selection restrictions identify this interaction. All expectations use a stated baseline population and probability law. Identification means every admissible data-generating model with the same observed-data law yields the same displayed target; otherwise report the identified set. Exact equations, horizons and assumptions are editorial, rather than source-given canonical conjectures.
\subsection*{Variants and assumption changes}
\begin{enumerate}
\item \textbf{Editorial total-policy variant.} Estimate \(E[Y_{imt}(1)-Y_{imt}(0)-Y_{ift}(1)+Y_{ift}(0)]\) on a preannouncement couple cohort allowing births to change. This includes fertility selection.
\item \textbf{Editorial mechanism variant.} Factorially vary childcare and father-only leave with support in all arms. Earnings interactions do not separately establish discrimination, tastes or bargaining without extra restrictions.
\end{enumerate}
\subsection*{References and source audit}
\textbf{1.} Existing family-policy evidence; exact child-penalty mechanism experiment not verified as source question. Locator: NBER Working Paper 30685, November 2022; abstract. Primary indexed abstract inspected; direct page failed. URL: \url{https://www.nber.org/papers/w30685}.

\textbf{2.} Existing cross-country child-penalty estimates; policy mechanism contrast is editorial. Locator: Review of Economic Studies 92(5), 3174–3207; published online 24 October 2024; abstract and Section 1. Full text or primary publication page inspected. URL: \url{https://academic.oup.com/restud/article/92/5/3174/7840285}.
\begin{enumerate}\raggedright
\item\hypertarget{definitions:source:30007739:1}{} Stefania Albanesi; Claudia Olivetti; Barbara Petrongolo. 2022. \emph{Families, Labor Markets, and Policy}. NBER Working Paper 30685, November 2022; abstract.
\url{https://www.nber.org/papers/w30685}.
DOI: \url{https://doi.org/10.3386/w30685}.
\item\hypertarget{definitions:source:30007739:2}{} Henrik Kleven; Camille Landais; Gabriel Leite-Mariante. 2025. \emph{The Child Penalty Atlas}. Review of Economic Studies 92(5), 3174–3207; published online 24 October 2024; abstract and Section 1.
\url{https://academic.oup.com/restud/article/92/5/3174/7840285}.
DOI: \url{https://doi.org/10.1093/restud/rdae104}.
\end{enumerate}

\subsection*{Common conventions: Source mathematical conventions}

These conventions apply to each entry unless explicitly replaced.

\textbf{Models and identification.} A primitive model \(m\) specifies all probability laws, feasible choices, information, equilibrium and measurement rules used by its target. The admissible class \(\mathcal M\) is fixed independently of outcomes. If \(P_O(m)\) is its observable probability law and \(\tau(m)\) the target, its identified set at observable law \(P\) is
\[
 \mathcal I_\tau(P)=\{\tau(m):m\in\mathcal M,\ P_O(m)=P\}.
\]
Point identification means that this set is a singleton. An empty set signals incompatibility of data and assumptions. Coordinate bounds need not describe the joint set. Bounds are sharp only if every retained target is attainable or arbitrarily approachable by compatible models. Finite bounds require explicit restrictions on unobserved quantities; unrestricted confounding, hidden wealth, extrapolation or arbitrary equilibrium selection can yield uninformative sets.

\textbf{Causal interventions.} A potential outcome \(Y_i(p)\) is the outcome under a complete policy regime \(p\), including timing, financing, information and market spillovers. The target population and units are fixed before treatment. With interference, use the complete assignment vector or a justified exposure mapping; individual no-interference notation alone cannot identify an economy-wide intervention. A difference of mean potential outcomes does not require the cross-world joint distribution. Conditioning on post-treatment migration, survival, enrollment or employment changes the estimand and needs separate justification.

\textbf{Statistical statements.} An estimator, confidence set, forecast or test specifies a sampling law, model class and loss/coverage criterion. Uniform \((1-\alpha)\)-coverage means \(\inf_{m\in\mathcal M}\Pr_m\{\tau(m)\in C_n\}\geq1-\alpha\), or an explicitly asymptotic version. The whole parameter space trivially has coverage, so informativeness requires a stated width, risk or power objective. A forecasting improvement is relative to a named benchmark, information set and evaluation population.

\textbf{Equilibrium and optimization.} An equilibrium comprises feasible plans, optimality, beliefs where needed, budget constraints and all market/resource clearing. The primitives or a stated benchmark define each correspondence \(\mathcal E(p;m)\). Nonemptiness is assumed or proved, never inferred just from compact policy sets. A selected-equilibrium target uses a declared selection rule; a robust target quantifies over all permitted equilibria. Use suprema or infima unless attainment is established. An \(\varepsilon\)-optimal policy is feasible and within \(\varepsilon\) of the finite optimum value. Report an empty feasible set separately.

\textbf{Welfare and accounting.} Normative utility, interpersonal normalization, social weights, numeraire, discounting and the use of public receipts are primitives. Behavioral utility need not coincide with the welfare criterion. Household welfare includes ownership income and rents once; adding profits again to compensating variation generally double-counts them. A monetary equivalent is defined by an explicit transfer and price convention, with monotonicity and range conditions. Average effects, mechanism contributions, transfers and welfare losses are different targets. Infinite sums require convergence and dynamic feasible plans require terminal or no-Ponzi conditions.

\textbf{Source versions.} A working-paper date, revision date and journal publication year are distinguished. A DOI for a journal version is not automatically the DOI of an earlier manuscript. Locators refer to the stated version and printed pages where specified. A metadata-only check does not verify a claim about a full-text conclusion. Every entry's source subsection narrows the provenance claim accordingly.
% END ECONOMICS STATEMENT
\end{document}
